\chapter{ASS1ST selection round (menu 23)}
\label{ch:menu23}
\menulabel{23}

\section{Purpose}

One selection step of ASS1ST: read the round's perturbation-theory density,
diagonalize its internal/internal and external/external blocks \emph{separately}
(the source's construction of quasi-natural orbitals), and propose the next
active space from the threshold band -- growing and shrinking both allowed.

\section{What you need}

The \code{orca\_2json} export of a finished round, with the request
\code{\{"MOCoefficients": true, "1elIntegrals": ["S"], "Densities": ["all"]\}}
(menu~\menu{22} writes it); the run must have produced its density sidecar
(\code{KeepDens} plus the \code{PTSettings} block, both in the generated
input).

\section{Walkthrough}

\lstinputlisting[style=fbkcode, firstline=54, lastline=111,
  title={The menu-23 example session (the N2 CAS(6,6) round at the relaxed
  0.03 band: the internal block is nearly diagonal, the external block is not,
  and the sigma pair drifts out -- the next space is the (4e, 4o) pi/pi*
  quartet, and the next round's input is written beside the report)}]
  {examples/expected/m23_ass1st_round.txt}

\section{What you get}

The block quasi-occupation tables with the band marked, the active orbitals
with their CASSCF occupations, the next-space suggestion, a report
(\file{<export>.ass1st.fbk.md}) with the citations, and -- unless the round is
self-consistent -- the next round's input \file{<stem>.r\{N+1\}.inp} with its
\code{.json.conf}.  Given an mkl template of the run (an \code{orca\_2mkl}
copy of its \code{.gbw}), the menu additionally writes the round's
quasi-natural orbital set as \file{<export>.qno.fbk.mkl}, ordered so that the
engine's by-orbital-order window takes the inactive prefix and then the
suggested active window; feeding it back as the next round's orbital guess
reproduces the suggestion directly.

\section{How to read it}

\begin{readbox}
\begin{itemize}
  \item every quasi-occupation is a \emph{block} eigenvalue of the exported
    NEVPT2 unrelaxed density.  ORCA's own printed ``Natural Orbital Occupation
    Numbers'' block is the whole-space naturalization and is \emph{not} what
    the scheme reads (measured); the external block in particular is far from
    diagonal, so the block step is what makes the external side meaningful;
  \item band bookkeeping (the source's): an internal quasi-orbital inside the
    band adds $+2$ electrons and $+1$ orbital to the active space, an external
    one adds $+0$ electrons and $+1$ orbital; an active orbital whose CASSCF
    occupation has drifted to the top of the band is reassigned to internal
    ($-2$~e, $-1$~o), at the bottom to external ($-0$~e, $-1$~o);
  \item \code{self-consistent} means no orbital crossed the band and no active
    orbital drifted: use that space for the production calculation.  If the
    same space was visited before, the source warns of cycling between spaces
    -- break the tie by chemical judgment;
  \item the density is ORCA's \textbf{FIC-NEVPT2 unrelaxed density}, while the
    source works with the SC-NEVPT2 first-order density; ORCA answers the
    unrelaxed density only for the FIC ansatz (measured: SC plus
    \code{Density Unrelaxed} is refused).  The report states this substitution
    with every round;
  \item the source's own caveats travel with the output: the result depends on
    the initial space and on the potential-energy point -- run the scheme at
    the geometries you care about.
\end{itemize}
\end{readbox}

\section{Measured chain behaviour}

On the N2/def2-SVP fixtures, round~1 CAS(6,6) suggests shrinking the sigma
pair out -- $(6,6) \to (4,4)$, the $\pi/\pi^{*}$ quartet -- and a rerun of that
space from the default guess converged to a \emph{different-shaped} $(4,4)$
solution, which the next analysis then reports as a further shrink.  This is
the multi-solution behaviour this project records elsewhere, and the reason
for the boundary note in Chapter~\ref{ch:menu22}: check the round's active
orbitals each time before continuing a chain.
